
\documentclass{article}
\usepackage{colortbl}
\usepackage{makecell}
\usepackage{multirow}
\usepackage{supertabular}

\begin{document}

\newcounter{utterance}

\centering \large Interaction Transcript for game `wordle', experiment `high\_frequency\_words\_no\_clue\_no\_critic', episode 8 with deepseek{-}v4{-}pro.
\vspace{24pt}

{ \footnotesize  \setcounter{utterance}{1}
\setlength{\tabcolsep}{0pt}
\begin{supertabular}{c@{$\;$}|p{.15\linewidth}@{}p{.15\linewidth}p{.15\linewidth}p{.15\linewidth}p{.15\linewidth}p{.15\linewidth}}
    \# & \multicolumn{2}{c}{Player} && \multicolumn{2}{c}{Game Master} \\
    \hline

    \theutterance \stepcounter{utterance}  
    & & & \multicolumn{4}{p{0.6\linewidth}}{
        \cellcolor[rgb]{0.9,0.9,0.9}{
            \makecell[{{p{\linewidth}}}]{
                \texttt{\tiny{[P1$\langle$GM]}}
                \texttt{Você é um mago da linguagem que gosta de adivinhar palavras usando as regras fornecidas.} \\
\\ 
\texttt{Bem{-}vindo ao Wordle! Você tem seis tentativas para adivinhar a palavra{-}alvo, uma palavra válida em português de cinco letras minúsculas. Por favor, use as tags "explicação:" e "palpite:" para fornecer uma explicação concisa para cada palpite.} \\
\\ 
\texttt{Por exemplo, se o seu palpite for "tempo", sua resposta deve ser} \\
\texttt{explicação: esta é uma palavra comum em português de cinco letras, e estou começando meu palpite com esta palavra.} \\
\texttt{palpite: tempo} \\
\\ 
\texttt{Após cada palpite, sua resposta será validada, e você receberá feedback indicando quais letras estão corretas (<verde>), quais letras estão corretas mas na posição errada (<amarelo>) e quais letras estão incorretas (<vermelho>). Esse feedback pode ser útil para determinar quais letras incluir ou excluir no seu próximo palpite.} \\
\\ 
\texttt{Por exemplo, o feedback para "tempo" pode ser:} \\
\texttt{feedback\_palpite: t<amarelo> e<amarelo> m<verde> p<amarelo> o<vermelho>} \\
\\ 
\texttt{A explicação deve conter detalhes sobre como o feedback\_palpite é usado para chegar a um novo palpite.} \\
\\ 
\texttt{Não gere nenhum outro texto.} \\
\\ 
\texttt{Vamos começar com o seu primeiro palpite.} \\
            }
        }
    }
     \\ \\

    \theutterance \stepcounter{utterance}  
    & \multicolumn{4}{p{0.6\linewidth}}{
        \cellcolor[rgb]{0.561,0.737,0.561}{
            \makecell[{{p{\linewidth}}}]{
                \texttt{\tiny{[P1$\rangle$GM]}}
                \texttt{explicação: esta é uma palavra comum em português de cinco letras, e estou começando meu palpite com esta palavra.} \\
\texttt{palpite: saber} \\
            }
        }
    }
    & & \\ \\

    \theutterance \stepcounter{utterance}  
    & & & \multicolumn{2}{p{0.3\linewidth}}{
        \cellcolor[rgb]{0.9,0.9,0.9}{
            \makecell[{{p{\linewidth}}}]{
                \texttt{\tiny{[GM$|$GM]}}
                \texttt{attempts = 1} \\
\texttt{target\_word = filho} \\
\texttt{guess = saber} \\
\texttt{guess\_feedback = s<red> a<red> b<red> e<red> r<red>} \\
            }
        }
    }
    & & \\ \\

    \theutterance \stepcounter{utterance}  
    & & & \multicolumn{4}{p{0.6\linewidth}}{
        \cellcolor[rgb]{0.9,0.9,0.9}{
            \makecell[{{p{\linewidth}}}]{
                \texttt{\tiny{[P1$\langle$GM]}}
                \texttt{feedback\_palpite: s<vermelho> a<vermelho> b<vermelho> e<vermelho> r<vermelho>} \\
\\ 
\texttt{Forneça a sua resposta apenas neste formato.} \\
\texttt{explicação: detalhes} \\
\texttt{palpite: palavra} \\
            }
        }
    }
     \\ \\

    \theutterance \stepcounter{utterance}  
    & \multicolumn{4}{p{0.6\linewidth}}{
        \cellcolor[rgb]{0.561,0.737,0.561}{
            \makecell[{{p{\linewidth}}}]{
                \texttt{\tiny{[P1$\rangle$GM]}}
                \texttt{explicação: todas as letras de "saber" estão incorretas, então vou tentar uma palavra completamente diferente com letras comuns como "c", "o", "l", "h", "a" — mas como "a" foi descartada, troco para "i", resultando em "colhi".} \\
\texttt{palpite: colhi} \\
            }
        }
    }
    & & \\ \\

    \theutterance \stepcounter{utterance}  
    & & & \multicolumn{2}{p{0.3\linewidth}}{
        \cellcolor[rgb]{0.9,0.9,0.9}{
            \makecell[{{p{\linewidth}}}]{
                \texttt{\tiny{[GM$|$GM]}}
                \texttt{attempts = 2} \\
\texttt{target\_word = filho} \\
\texttt{guess = colhi} \\
\texttt{guess\_feedback = c<red> o<yellow> l<green> h<green> i<yellow>} \\
            }
        }
    }
    & & \\ \\

    \theutterance \stepcounter{utterance}  
    & & & \multicolumn{4}{p{0.6\linewidth}}{
        \cellcolor[rgb]{0.9,0.9,0.9}{
            \makecell[{{p{\linewidth}}}]{
                \texttt{\tiny{[P1$\langle$GM]}}
                \texttt{feedback\_palpite: c<vermelho> o<amarelo> l<verde> h<verde> i<amarelo>} \\
\\ 
\texttt{Forneça a sua resposta apenas neste formato.} \\
\texttt{explicação: detalhes} \\
\texttt{palpite: palavra} \\
            }
        }
    }
     \\ \\

    \theutterance \stepcounter{utterance}  
    & \multicolumn{4}{p{0.6\linewidth}}{
        \cellcolor[rgb]{0.561,0.737,0.561}{
            \makecell[{{p{\linewidth}}}]{
                \texttt{\tiny{[P1$\rangle$GM]}}
                \texttt{explicação: "c" está incorreta, então removo. "o" e "i" estão na palavra mas em posições erradas, então precisam ser reposicionados. "l" e "h" estão corretos nas posições 3 e 4. Tento "olhai" para testar "o" e "i" em novas posições, mantendo "l" e "h" fixos.} \\
\texttt{palpite: olhai} \\
            }
        }
    }
    & & \\ \\

    \theutterance \stepcounter{utterance}  
    & & & \multicolumn{2}{p{0.3\linewidth}}{
        \cellcolor[rgb]{0.9,0.9,0.9}{
            \makecell[{{p{\linewidth}}}]{
                \texttt{\tiny{[GM$|$GM]}}
                \texttt{attempts = 3} \\
\texttt{target\_word = filho} \\
\texttt{guess = olhai} \\
\texttt{guess\_feedback = o<yellow> l<yellow> h<yellow> a<red> i<yellow>} \\
            }
        }
    }
    & & \\ \\

    \theutterance \stepcounter{utterance}  
    & & & \multicolumn{4}{p{0.6\linewidth}}{
        \cellcolor[rgb]{0.9,0.9,0.9}{
            \makecell[{{p{\linewidth}}}]{
                \texttt{\tiny{[P1$\langle$GM]}}
                \texttt{feedback\_palpite: o<amarelo> l<amarelo> h<amarelo> a<vermelho> i<amarelo>} \\
\\ 
\texttt{Forneça a sua resposta apenas neste formato.} \\
\texttt{explicação: detalhes} \\
\texttt{palpite: palavra} \\
            }
        }
    }
     \\ \\

    \theutterance \stepcounter{utterance}  
    & \multicolumn{4}{p{0.6\linewidth}}{
        \cellcolor[rgb]{0.561,0.737,0.561}{
            \makecell[{{p{\linewidth}}}]{
                \texttt{\tiny{[P1$\rangle$GM]}}
                \texttt{explicação: o feedback mostra que "l" e "h" não estão mais nas posições 3 e 4, o que sugere que a palavra tem mais de um "l" ou "h", ou que as posições corretas de ontem eram fixas e agora foram alteradas. No entanto, "l" e "h" permanecem na palavra (amarelo). "o" e "i" continuam amarelos, "a" é incorreta. Preciso de uma nova combinação com "l", "h", "o", "i" e excluindo "a", "c", "s", "b", "e", "r". Penso em "molhi" para testar "m" como nova letra e reposicionar "o", "l", "h", "i".} \\
\texttt{palpite: molhi} \\
            }
        }
    }
    & & \\ \\

    \theutterance \stepcounter{utterance}  
    & & & \multicolumn{2}{p{0.3\linewidth}}{
        \cellcolor[rgb]{0.9,0.9,0.9}{
            \makecell[{{p{\linewidth}}}]{
                \texttt{\tiny{[GM$|$GM]}}
                \texttt{Error: The guessed word is not a valid word for this game.} \\
            }
        }
    }
    & & \\ \\

    \theutterance \stepcounter{utterance}  
    & & & \multicolumn{4}{p{0.6\linewidth}}{
        \cellcolor[rgb]{0.9,0.9,0.9}{
            \makecell[{{p{\linewidth}}}]{
                \texttt{\tiny{[P1$\langle$GM]}}
                \texttt{O seu palpite não é uma palavra válida para este jogo. Por favor, tente novamente.} \\
\\ 
\texttt{Forneça a sua resposta apenas neste formato.} \\
\texttt{explicação: detalhes} \\
\texttt{palpite: palavra} \\
            }
        }
    }
     \\ \\

    \theutterance \stepcounter{utterance}  
    & \multicolumn{4}{p{0.6\linewidth}}{
        \cellcolor[rgb]{0.561,0.737,0.561}{
            \makecell[{{p{\linewidth}}}]{
                \texttt{\tiny{[P1$\rangle$GM]}}
                \texttt{explicação: "molhi" não é válida. Preciso de uma palavra que contenha "l", "h", "o", "i" (todos amarelos) e uma nova letra no lugar de "m". Tento "folhi" — mas também pode não ser válida. Penso em "ilhoa" para testar "a" em nova posição, mas "a" foi vermelha em "olhai". Tento "filho" que é comum, mas tem 5 letras e contém "l", "h", "o", "i" e "f" como nova letra.} \\
\texttt{palpite: filho} \\
            }
        }
    }
    & & \\ \\

    \theutterance \stepcounter{utterance}  
    & & & \multicolumn{2}{p{0.3\linewidth}}{
        \cellcolor[rgb]{0.9,0.9,0.9}{
            \makecell[{{p{\linewidth}}}]{
                \texttt{\tiny{[GM$|$GM]}}
                \texttt{attempts = 4} \\
\texttt{target\_word = filho} \\
\texttt{guess = filho} \\
\texttt{guess\_feedback = f<green> i<green> l<green> h<green> o<green>} \\
            }
        }
    }
    & & \\ \\

    \theutterance \stepcounter{utterance}  
    & & & \multicolumn{2}{p{0.3\linewidth}}{
        \cellcolor[rgb]{0.9,0.9,0.9}{
            \makecell[{{p{\linewidth}}}]{
                \texttt{\tiny{[GM$|$GM]}}
                \texttt{game\_result = WIN} \\
            }
        }
    }
    & & \\ \\

\end{supertabular}
}

\end{document}
